This part is about the brains of our agents, and how they use them to learn.  
The previous part showed us how a basic next-token LLM is trained. In agentic models typically  more advanced reasoning and multi-modal models are used (which we continue to call {\em LLM}).  In this part we will discuss these more advanced models, why they are needed, and how they are trained and used.  

The first chapter of this part introduces reasoning models. 
Originally LLMs could not solve simple reasoning problems, and the chain of thought concept caused a major leap forward in this field. We will discuss math word problems, chain of thought prompting, and in-context learning. We will see models that use self-reflection to try to improve their own behavior. Self-reflection loops led to the development of agentic harnesses for LLMs. 

In the following chapter we go beyond language: we will see that the transformer architecture also works for vision and action domains. We will discuss various kinds of World Models, models  that allow agents to act more intelligently in their interaction with the world. 

In the last  chapter of this part, we introduce the concept of metalearning, a paradigm that is present in many of the types of learning that we see in agentic LLMs. We will discuss what metalearning is, and why it is a challenging form of machine learning.
